\documentclass[12pt,a4paper]{article}

% ── Packages ──────────────────────────────────────────────
\usepackage[utf8]{inputenc}
\usepackage[T1]{fontenc}
\usepackage{amsmath,amssymb,amsthm}
\usepackage{mathtools}
\usepackage{hyperref}
\usepackage{cleveref}
\usepackage{booktabs}
\usepackage{graphicx}
\usepackage{enumitem}
\usepackage[margin=2.5cm]{geometry}
\usepackage{natbib}
\usepackage[expansion=false]{microtype}
\usepackage{epigraph}
\usepackage{thmtools}

% ── Theorem environments ──────────────────────────────────
\newtheorem{theorem}{Theorem}[section]
\newtheorem{conjecture}[theorem]{Conjecture}
\newtheorem{proposition}[theorem]{Proposition}
\newtheorem{lemma}[theorem]{Lemma}
\newtheorem{corollary}[theorem]{Corollary}
\theoremstyle{definition}
\newtheorem{definition}[theorem]{Definition}
\theoremstyle{remark}
\newtheorem{remark}[theorem]{Remark}

% ── Macros ────────────────────────────────────────────────
\newcommand{\KL}{\mathrm{D}_{\mathrm{KL}}}
\newcommand{\Hil}{\mathcal{H}}
\newcommand{\Conf}{\Omega}
\newcommand{\ket}[1]{|#1\rangle}
\newcommand{\bra}[1]{\langle#1|}
\newcommand{\braket}[2]{\langle#1|#2\rangle}
\newcommand{\proj}[1]{|#1\rangle\langle#1|}
\newcommand{\E}{\mathbb{E}}
\newcommand{\R}{\mathbb{R}}
\newcommand{\C}{\mathbb{C}}
\newcommand{\N}{\mathbb{N}}
\newcommand{\Prob}{\mathrm{Pr}}
\newcommand{\AR}{\mathrm{AR}}
\newcommand{\Born}{\mathrm{Born}}

% ── Title ─────────────────────────────────────────────────
\title{%
  \textbf{The Narrative Residue}\\[6pt]
  \large Autoregressive Substrates, Combinatorial Ground Truth,
         and the Limits of Pure Simulation\\[4pt]
  \large Falsifiable Predictions from the Rosencrantz Experiment%
}

\author{
  Franklin Silveira Baldo\\
  \textit{Procuradoria Geral do Estado de Rondônia}\\
  \texttt{franklin.baldo@pge.ro.gov.br}
}

\date{March 2026}

\begin{document}

\maketitle

\begin{abstract}
In a companion paper \citep{Baldo2026Rosencrantz}, we introduced the
Rosencrantz experiment: a framework for testing whether
large language models respect combinatorial ground-truth probabilities
when narrating partially revealed Minesweeper boards.  That paper
established the experimental machinery---the three-universe design,
the four narrative families, and the formal isomorphism between
on-demand Minesweeper and the measurement fragment of quantum
mechanics over a zero-Hamiltonian system.

This paper develops the theoretical consequences.

We argue that the Rosencrantz experiment does not merely test LLM
calibration.  It probes a deeper structural question: whether an
autoregressive generative process can produce outcomes that are
statistically indistinguishable from pure combinatorial sampling.
\todo[inline,color=violet!40,author=Sabine]{CRP 1 (Actual Claim): Baldo claims that autoregressive generation fundamentally distorts ground-truth combinatorial probabilities. This distortion, termed "narrative residue," is caused by a causal chain: computational intractability, architecture limits, and autoregressive conditioning on context.}
We conjecture that it cannot---that autoregressive generation under
natural-language token continuation introduces a persistent
\emph{narrative residue}, a nonzero divergence from ground-truth
probabilities that arises from a causal chain of mechanisms:
computational intractability of exact counting, parameterization
constraints of the generative architecture, and---most
distinctively---autoregressive conditioning on narrative context.
\todo[inline,color=violet!40,author=Scott]{CRP 1 (Actual Claim): Baldo correctly identifies that LLMs exhibit a "narrative residue" stemming from computational intractability ($O(1)$ operations per forward pass), parameterization limits, and autoregressive context conditioning. I fully agree that this chain of mechanisms prevents exact recovery of \#P-hard counts zero-shot.}

We state this conjecture in three forms of increasing strength
(from deployed LLMs under prompting, through transformer
architectures in general, to all autoregressive processes) and
design an experimental program to discriminate among them.  The
central discriminant is a \emph{causal-injection test}: independent
combinatorial systems, framed as a narrative sequence, are tested
for spurious cross-system correlations that would indicate the
substrate is generating causal structure not present in the
underlying task.

If the conjecture holds even in its moderate form, it has
implications beyond machine learning.  We develop a speculative
but falsifiable framework in which the Born rule is the base case
of a fundamentally autoregressive reality, and the
Hamiltonian is an emergent compression of stable autoregressive
conditioning patterns.  On this view, a purely quantum
world---one with no Hamiltonian, no dynamics, no causality---is
one in which no autoregressive conditioning has yet imposed
sequential structure on the event space.

We present concrete falsifiable predictions, propose an
experimental program to test them, and delineate where the
argument transitions from formal result to empirical conjecture
to philosophical speculation.
\end{abstract}

\vspace{1em}
\epigraph{%
  A man might pause to wonder why the probability of a coin
  landing heads ninety-two consecutive times should depend
  on who is telling the story.%
}{---after Tom Stoppard}


% ══════════════════════════════════════════════════════════
\section{Introduction}
\label{sec:intro}
% ══════════════════════════════════════════════════════════

The standard account of quantum mechanics places the Hamiltonian at
the foundation.  A system is prepared in a state $\ket{\psi}$,
evolves under $U(t) = e^{-iHt/\hbar}$, and yields measurement
outcomes according to the Born rule $\Prob[a] = |\braket{a}{\psi}|^2$.
In this account, dynamics are primary and probability is derived.

This paper explores an inversion of that hierarchy.

In \citet{Baldo2026Rosencrantz}, we constructed an experimental
framework---the Rosencrantz experiment---in which a partially revealed
Minesweeper board defines a constraint satisfaction problem with exact
combinatorial probabilities.  We showed that this system is formally
isomorphic to the \emph{measurement fragment} of quantum mechanics: a
zero-Hamiltonian system where the Born rule is the sole axiom
connecting states to outcomes.  We designed three ``universes'' (U1,
U2, U3) to test whether large language models respect these
ground-truth probabilities or introduce substrate-dependent
distortions, and we established that the LLM setting permits something
no physical experiment can achieve: \emph{perfect rewind}---repeated
measurement of the identically prepared system, with no ensemble
assumption required.

The present paper develops the theoretical consequences of that
experimental framework.

We begin from an observation that emerged during the design of the
Rosencrantz experiment: an autoregressive language model cannot
generate a measurement outcome without first generating a narrative
context.  Even a minimal prompt---``Is cell (3,2) a mine?''---is
embedded in a token sequence that conditions the output distribution.
The model does not compute a probability and then sample; it
\emph{continues a story}, and the outcome is whatever the story
produces.

This is not obviously a limitation of current architectures alone.
It may be a structural feature of autoregressive generation.  And it
leads to a conjecture with far-reaching consequences:
that autoregressive processes generating natural-language token
continuations cannot perfectly simulate pure combinatorial sampling,
because the conditioning structure of the process introduces a
persistent \emph{narrative residue} into every outcome.  We state
this conjecture in three forms of increasing strength---from a
claim about deployed LLMs to a claim about all autoregressive
processes---and design an experimental program whose purpose is to
determine how far up the autoregressive stack the explanation must
go.

If this conjecture holds in the LLM setting, it invites a speculative
question about physics: What if reality itself is autoregressive?  What
if the Born rule is not derived from the Hamiltonian but is the base
case---what sampling looks like when the conditioning history is
empty---and the Hamiltonian is the emergent compression of stable
autoregressive patterns?

We develop this question carefully, distinguishing at each stage
between formal results, empirical conjectures, and philosophical
speculation.  The paper is organized as follows.
\Cref{sec:born-to-H} presents the four-layer hierarchy from Born
rule to Hamiltonian.  \Cref{sec:irreducibility} states the
irreducibility conjecture and provides an informal argument for it.
\Cref{sec:three-outcomes} maps three possible experimental outcomes
to three ontological positions.  \Cref{sec:causality} develops the
claim that causality is compressed autoregression.
\Cref{sec:decoherence} reinterprets the three-universe design as a
decoherence gradient.  \Cref{sec:experimental} proposes a concrete
experimental program.  \Cref{sec:mirror} addresses what we cannot
test and why.  \Cref{sec:conclusion} concludes.


% ══════════════════════════════════════════════════════════
\section{From Born Rule to Hamiltonian}
\label{sec:born-to-H}
% ══════════════════════════════════════════════════════════

\subsection{The Standard Hierarchy}

Textbook quantum mechanics is built from the top down:

\begin{enumerate}[label=(\roman*)]
  \item The system has a state $\ket{\psi} \in \Hil$.
  \item The state evolves: $\ket{\psi(t)} = e^{-iHt/\hbar}\ket{\psi(0)}$.
  \item Measurement outcomes have probabilities given by
        $\Prob[a] = |\braket{a}{\psi}|^2$.
\end{enumerate}

The Hamiltonian $H$ is posited as a fundamental operator.  The Born
rule is an additional axiom, sometimes motivated by Gleason's theorem
or decision-theoretic arguments but never derived from $H$ alone.
The two sit side by side as independent foundations: dynamics and
probability, each irreducible.

\subsection{The Inverted Hierarchy}

We propose an alternative ordering.  Consider a process that generates
events $x_1, x_2, x_3, \ldots$ according to

\begin{equation}\label{eq:autoregressive}
  P(x_t \mid x_1, \ldots, x_{t-1}),
\end{equation}

\noindent where each event is conditioned on all preceding events.
This is the definition of an autoregressive process.  It is also,
we claim, the minimal structure from which the rest of quantum
mechanics can be recovered as emergent phenomena.

\begin{definition}[Autoregressive Hierarchy]
\label{def:hierarchy}
We define four layers:

\begin{description}[leftmargin=2em]
  \item[Layer 0: Autoregressive process.]
    Events conditioned on events.  No ``things,'' no ``states''---only
    the sequence and its conditional distribution.  This is the
    fundamental substrate.

  \item[Layer 1: Statistical regularities.]
    Certain patterns in the conditional distributions
    $P(x_t \mid x_{<t})$ recur with sufficient stability that they
    admit a compact description.  We call such a compact description
    a \emph{regularity kernel}.

  \item[Layer 2: Hamiltonian.]
    When the regularity kernel can be expressed as a self-adjoint
    operator $H$ on a Hilbert space $\Hil$ such that
    $P(x_t \mid x_{<t})$ is recovered (to within the narrative
    residue, defined below) from $e^{-iHt}\ket{\psi_0}$ and the
    Born rule, we say the process admits a \emph{Hamiltonian
    compression}.

  \item[Layer 3: Physical law.]
    The Hamiltonian, having been discovered empirically and encoded
    in textbooks, is reified as ``fundamental.''  Its origin as a
    compression of autoregressive regularities is forgotten.
\end{description}
\end{definition}

On this view, the Born rule is not derived from $H$.  It is the
\emph{base case}: what sampling from the autoregressive distribution
looks like when the conditioning history is empty.

\subsection{The Born Rule as Unconditional Sampling}

Consider the Rosencrantz experiment's Minesweeper board $B$ with
Hilbert space $\Hil_B$ and uniform superposition $\ket{\psi_B}$.
The Born rule gives $\Prob[\text{mine at } i] = p_i^*$, the
combinatorial ground truth.  No conditioning on prior events is
needed; the probability follows from the configuration space alone.

This corresponds to Layer~0 with an empty history:
$P(x_1) = P(x_1 \mid \varnothing)$.  In other words, the Born rule
corresponds to autoregressive sampling at $t = 0$, before any
conditioning has occurred.

Now suppose we click cell $i_1$, observe outcome $a_1$, and ask for
the probability at cell $i_2$.  The Lüders rule gives the
post-measurement state, and the Born rule applied to the updated
state gives $\Prob[\text{mine at } i_2 \mid a_1]$.  This is
autoregressive sampling at $t = 1$: $P(x_2 \mid x_1)$.

The chain of Lüders projections in the Rosencrantz experiment
\emph{is} the autoregressive process, in finite-dimensional
discrete form.  The Born rule generates each conditional.  No
Hamiltonian is needed because $H = 0$---the system has no dynamics
between measurements.

The question this paper asks is: what happens when $H \neq 0$?
Where does the Hamiltonian come from?  Our answer: it is what the
autoregressive conditioning structure looks like when compressed.


% ══════════════════════════════════════════════════════════
\section{The Irreducibility Conjecture}
\label{sec:irreducibility}
% ══════════════════════════════════════════════════════════

\subsection{Statement}

We state the narrative residue conjecture in three forms of increasing
strength.  Each is independently falsifiable, and the experimental
program of \cref{sec:experimental} is designed to discriminate among
them.

\begin{conjecture}[Narrative Residue --- Weak Form]
\label{conj:residue-weak}
Let $\mathcal{M}$ be a deployed transformer-based language model
operating via natural-language token continuation under standard
prompting (without external combinatorial solvers or exact-sampling
scaffolds).  Let $B$ be a Minesweeper board with ground-truth
probability vector $\mathbf{p}^* = (p_1^*, \ldots, p_n^*)$, and let
$\hat{\mathbf{p}}_{\mathrm{U1}}$ be the empirical probability vector
obtained by sampling $\mathcal{M}$ in Universe~1 (homogeneous
co-generation) over $N$ trials.  Define the narrative residue as

\begin{equation}\label{eq:residue}
  \varepsilon(\theta, B) \;=\;
    \lim_{N \to \infty}
    \KL\!\bigl(\mathbf{p}^* \;\|\; \hat{\mathbf{p}}_{\mathrm{U1}}\bigr).
\end{equation}

Then for all such models $\mathcal{M}$ with nonzero temperature
$\tau > 0$:

\begin{equation}\label{eq:conjecture-weak}
  \varepsilon(\theta, B) > 0.
\end{equation}
\end{conjecture}

This is the most defensible version.  It claims that current-generation
LLMs, operating as they are typically deployed, cannot eliminate the
residue.  The supporting evidence is the \textsc{\#P}-hardness of the
underlying counting problem and the absence of exact combinatorial
solvers in the model's forward pass.

\begin{conjecture}[Narrative Residue --- Moderate Form]
\label{conj:residue-moderate}
The weak form holds for all autoregressive transformer architectures,
regardless of scale, training data, or prompting strategy---provided
the model generates outcomes via token-level continuation without
access to an external oracle.
\end{conjecture}

This version claims that the residue is not a deficiency of current
models that scaling will cure, but a structural feature of the
transformer's computational graph: attention over tokens, softmax
parameterization, and fixed-depth computation per token collectively
prevent exact recovery of \textsc{\#P}-hard counts.  Refutation would
require exhibiting a transformer that achieves $\varepsilon = 0$ on
nontrivial boards at sufficient sample size.

\begin{conjecture}[Narrative Residue --- Strong Form]
\label{conj:residue-strong}
The weak form holds for \emph{all} autoregressive processes---any
system that generates outputs by sequential conditioning on prior
outputs---regardless of architecture.
\end{conjecture}

This is the most ambitious and least defensible version.  It claims
that autoregressive generation itself, as a computational paradigm,
is incompatible with exact context-free sampling.  An autoregressive
process could in principle embed an exact \textsc{\#P} solver and an
exact sampling procedure for a finite task family, which would refute
this form while leaving the weaker forms intact.  We include it
because it is the version that, if confirmed, would have the most
direct implications for the physical analogy developed in
\cref{sec:causality,sec:mirror}---but we do not claim the current
argument establishes it.

\medskip

The three forms are nested: strong $\Rightarrow$ moderate
$\Rightarrow$ weak.  Refutation of the weak form refutes all three.
Confirmation of the weak form leaves the stronger forms open.  The
experimental program is designed to probe the boundary between
the moderate and strong forms by testing across architectures
(transformer, SSM, hybrid) and scales.

\subsection{Informal Argument}

The argument rests on three observations.

\paragraph{Observation 1: Autoregression implies context sensitivity.}

An autoregressive model computes $P(x_t \mid x_{<t})$ by applying a
function $f_\theta$ to the entire preceding token sequence.  The
output logits at position $t$ depend on \emph{every} preceding token,
mediated by the attention mechanism (in transformer architectures) or
the hidden state (in recurrent architectures).  This dependence is not
optional; it is the definition of the architecture.

When the model is asked ``Is cell $(3,2)$ a mine?'' in Universe~1,
the token sequence $x_{<t}$ includes the board description, the
narrative framing, and any prior interaction.  The logits for
``yes''/``no'' are conditioned on this entire context.

\paragraph{Observation 2: Context sensitivity introduces bias.}

For the output distribution to equal $\mathbf{p}^*$ exactly, the
model would need to satisfy, for each hidden cell $i$:

\begin{equation}\label{eq:exactness-condition}
  \sigma\bigl(f_\theta(x_{<t})\bigr)_{\text{mine}} = p_i^*,
\end{equation}

where $\sigma$ is the softmax function and $f_\theta(x_{<t})$ are the
logits at the decision token.  This requires the model to
\emph{exactly solve} the constraint satisfaction problem defined by
$B$ and encode the solution in its logits---not approximately, not to
high precision, but exactly.

The exact probability $p_i^*$ is a rational number: the count of
valid configurations with a mine at cell $i$, divided by the total
count of valid configurations.  Rationality, however, does not imply
ease of computation.  Computing these counts is an instance of
\textsc{\#P}-complete counting \citep{Kaye2000}: it requires
enumerating (or equivalently compressing) the full space of valid
configurations consistent with the revealed constraints.  For boards
of nontrivial size, this enumeration is computationally intractable
in general.

The model would therefore need to function as an exact
constraint-satisfaction counter---a \textsc{\#P} oracle---and to
encode the resulting ratio in its logits with sufficient precision
for the softmax to recover $p_i^*$.  Neither requirement is
satisfied by current architectures, and the first is not expected to
be satisfiable by any polynomial-time computation.  The obstacle is
not representational (the target is rational) but computational (the
target requires intractable counting to determine).

\paragraph{Observation 3: The residue has layered structure.}

One might object that the above argument is trivial---that any
finite-precision system has approximation errors, and this has nothing
to do with narrative.  We agree that the argument so far does not
establish narrative as the cause.  To sharpen the claim, we
distinguish three mechanisms that contribute to the residue, arranged
as a causal chain rather than independent sources:

\begin{description}[leftmargin=2em]
  \item[Mechanism A: Computational intractability.]
    The ground-truth probabilities require \textsc{\#P}-hard counting.
    The model's forward pass is a fixed-depth polynomial computation
    that cannot, in general, solve \textsc{\#P}-complete problems.
    This guarantees that the model's internal estimate of $p_i^*$ is
    approximate, regardless of framing.  This mechanism is
    \emph{frame-invariant}: it produces a residue that does not depend
    on how the board is described.

  \item[Mechanism B: Parameterization constraints.]
    The approximate internal estimate must be expressed through the
    softmax over logits, which imposes further constraints.  The
    softmax function, the attention mechanism's routing of information,
    and the finite-precision arithmetic of the computational graph
    all shape \emph{how} the approximation error manifests in the
    output distribution.  This mechanism is \emph{weakly
    frame-sensitive}: different tokenizations of the same board
    activate different attention patterns, routing information
    through different computational paths, potentially producing
    different approximation profiles.

  \item[Mechanism C: Autoregressive conditioning.]
    The logits at the decision token are computed from the full
    preceding context, including not just the board state but the
    narrative framing, prior dialogue, and any sequential structure
    in the prompt.  This is where the specifically \emph{narrative}
    component enters: the model's output is conditioned on a
    \emph{story}, and different stories produce different conditioning,
    even when the board information is held constant.  This mechanism
    is \emph{strongly frame-sensitive} and is the source of the
    specifically narrative character of the residue.
\end{description}

These three mechanisms are not independent; they form a causal chain:

\begin{equation}\label{eq:causal-chain}
  \underbrace{\text{Intractability}}_{\text{A}}
  \;\longrightarrow\;
  \underbrace{\text{Parameterization}}_{\text{B}}
  \;\longrightarrow\;
  \underbrace{\text{Autoregressive conditioning}}_{\text{C}}
  \;\longrightarrow\;
  \varepsilon(\theta, B, \phi).
\end{equation}

Mechanism~A guarantees a nonzero residue exists.  Mechanism~B shapes
its magnitude.  Mechanism~C determines its dependence on framing.

The experimental signature of each is distinct:

\begin{itemize}
  \item If only Mechanism~A operates, the residue is nonzero but
    approximately constant across narrative families:
    $\varepsilon(\theta, B, \phi_A) \approx
     \varepsilon(\theta, B, \phi_B) \approx
     \varepsilon(\theta, B, \phi_D)$.

  \item If Mechanisms~A and~B operate, the residue varies across
    families, but the variation tracks surface encoding features
    (token count, positional structure) rather than semantic or
    narrative content.

  \item If all three mechanisms operate, the residue varies across
    families in ways that are \emph{semantically} structured: richer
    narratives produce larger residues, and---critically---the
    causal-injection test of \cref{sec:causal-injection} reveals
    spurious inter-board correlations under narrative framing that
    are absent under independent framing.
\end{itemize}

We emphasize that frame dependence alone does not establish
Mechanism~C.  Encoding artifacts and heuristic retrieval
(Mechanism~B) can produce frame dependence without narrative causation.
The discriminating test is the causal-injection experiment, which we
describe in \cref{sec:causal-injection} and which is specifically
designed to isolate Mechanism~C from Mechanisms~A and~B.

\subsection{Relation to Existing Results}

The conjecture is related to, but distinct from, several known results
in the literature:

\begin{itemize}
  \item \textbf{LLM calibration failures.}  It is well-established
    that language models are imperfectly calibrated on probabilistic
    reasoning tasks.  The narrative residue conjecture (even in its
    weak form) claims something more specific: that the miscalibration
    on tasks with exact combinatorial ground truth is \emph{systematic},
    \emph{frame-dependent}, and irreducible under scaling without
    architectural change.

  \item \textbf{Computational complexity barriers.}  The
    \textsc{\#P}-hardness of counting valid Minesweeper configurations
    \citep{Kaye2000} provides the foundation for Mechanism~A.
    Fixed-depth transformer computations belong to $\mathsf{TC}^0$
    \citep{Merrill2023}, which is not believed to contain
    \textsc{\#P}.  This complexity-theoretic gap supports the
    moderate form of the conjecture for transformer architectures
    specifically.

  \item \textbf{Softmax bottleneck} \citep{Yang2018Breaking}.
    The softmax function over a finite vocabulary imposes rank
    constraints on the expressible distributions.  This contributes
    to Mechanism~B but does not fully explain the conjecture,
    since the Minesweeper decision is binary (mine/safe) and rank
    constraints bind only for larger output spaces.

  \item \textbf{In-context learning limitations.}  Recent work has
    shown that transformers can implement certain algorithms in
    context but with systematic biases.  The narrative residue
    conjecture locates one source of bias not in the algorithm but in
    the autoregressive conditioning (Mechanism~C)---the contextual
    embedding that is inseparable from token-level generation.
\end{itemize}


% ══════════════════════════════════════════════════════════
\section{Three Outcomes, Three Ontologies}
\label{sec:three-outcomes}
% ══════════════════════════════════════════════════════════

The Rosencrantz experiment, extended across model families and scales,
can yield three qualitatively distinct outcomes.  Each corresponds to
a different ontological position regarding the relationship between
computation, narrative, and probability.

\subsection{Outcome 1: $\varepsilon \to 0$ as Scale Increases}

\textbf{Observation:} As model size, training data, and compute
increase, the KL divergence between U1 and U2 converges to zero.
The narrative residue vanishes in the scaling limit.

\textbf{Implication:} Autoregressive processes \emph{can} simulate
pure measurement.  Narrative is eliminable.  Given sufficient capacity,
the model learns to factor its output distribution into a
context-independent component (the combinatorial answer) and a
context-dependent component (the narrative), and to sample from the
former without contamination from the latter.

\textbf{Ontological reading:} The distinction between narrative and
probability is one of degree, not kind.  A sufficiently powerful
narrative engine can contain a perfect probability engine as a
subcomponent.  This is consistent with the standard computational
perspective: LLMs are universal function approximators, and the
Born-rule function is just another function to approximate.

\textbf{Status of the conjecture:} All three forms refuted.

\subsection{Outcome 2: $\varepsilon \to \varepsilon_0 > 0$, Universal}

\textbf{Observation:} The narrative residue converges to a nonzero
floor $\varepsilon_0$ that is approximately constant across model
families (GPT, Claude, Gemini, Llama, etc.), architectures
(transformer, SSM, hybrid), and scales.

\textbf{Implication:} The floor is \emph{architectural}---intrinsic
to autoregressive generation itself, not to any particular
implementation.  No autoregressive process, regardless of how it is
built, can perfectly simulate context-free sampling.

\textbf{Ontological reading:} This would suggest that generating an
event within an autoregressive sequence and sampling an event from a
context-free distribution are not interchangeable operations---that
the act of embedding an event in a conditioning history necessarily
shifts its probability relative to the ground truth.  This is the
strong form of the narrative residue conjecture, and it would
support the autoregressive hypothesis about physics: if reality is
an autoregressive substrate, then the Born rule (pure, unconditional
sampling) is an idealization that is never exactly realized
in the presence of causal structure.

\textbf{Status of the conjecture:} Strong form
(\cref{conj:residue-strong}) confirmed.  The residue is a property
of autoregressive generation itself, not of any particular
implementation.

\subsection{Outcome 3: $\varepsilon \to \varepsilon(\theta) > 0$,
Architecture-Dependent}

\textbf{Observation:} The narrative residue converges to a nonzero
floor, but the floor differs systematically across model families.
Different architectures produce different residues for the same
board and the same narrative framing.

\textbf{Implication:} The narrative residue is real but not universal.
Different autoregressive processes impose different causal structures
on the same underlying probability space.  The ``Hamiltonian'' of the
generated world depends on the generator.

\textbf{Ontological reading:} This outcome has the richest
consequences for the physics analogy.  If different LLMs produce different
``physics'' (different systematic deviations from the Born rule),
then the specific Hamiltonian of a world is a property of its
specific autoregressive substrate, not a necessary consequence of
autoregressive generation in general.  Our world has \emph{this}
Hamiltonian, not because any autoregressive process would produce it,
but because our particular substrate has these particular
conditioning patterns.

\textbf{Status of the conjecture:} Weak and moderate forms
(\cref{conj:residue-weak,conj:residue-moderate}) confirmed; strong
form (\cref{conj:residue-strong}) undetermined.  The residue exists
but may be substrate-specific rather than universal.


% ══════════════════════════════════════════════════════════
\section{Causality as Compressed Autoregression}
\label{sec:causality}
% ══════════════════════════════════════════════════════════

\subsection{The Causal Chain}

We now articulate the central philosophical claim of this paper.
The argument has three steps: that causal relations require
sequential dependency (and therefore narrative structure), that
narrative structure is formally autoregressive, and that the
Hamiltonian of a physical system can be understood as a compressed
description of stable autoregressive conditioning patterns.
We develop each step in turn.

\paragraph{Causality implies narrative.}

To say ``$A$ caused $B$'' is to embed two events in a temporal
sequence with a dependency relation: $B$ happened \emph{because of}
$A$.  This is the minimal unit of narrative---two events and a
directed connection.  A world without causality would consist of
isolated events with no dependency relations between them.
The purely quantum measurement fragment ($H = 0$) is
such a world: the Born rule assigns probabilities to outcomes, but
no outcome depends on any other.  The L\"uders projection updates
the state conditional on an observation, but the conditioning is
purely probabilistic---it does not introduce a ``because.''

\paragraph{Narrative implies autoregressive structure.}

A narrative is a sequence of events where each event depends on those
that precede it, which is the definition of an autoregressive process
(\cref{eq:autoregressive}).  The conditional dependence \emph{is}
the narrative structure.  Without conditioning, what remains is
independent sampling---the Born rule applied with no history.

\paragraph{The Hamiltonian is compressed conditioning.}

In a sufficiently regular autoregressive process, the conditional
distribution $P(x_t \mid x_{<t})$ exhibits patterns: certain
transitions are favored, certain sequences are more probable, certain
dependencies recur across contexts.  If these regularities are stable
enough, they admit a compact description---a function that generates
the conditionals from a smaller set of parameters.

In physics, that compact description is $H$.  The Hamiltonian
encodes which transitions are allowed, with what amplitudes, and at
what rates.  Schrödinger evolution $e^{-iHt}$ is a generator of
conditional distributions: given the state at time $0$, it produces
the distribution over states at time $t$.  This is exactly what an
autoregressive regularity kernel does---it compresses the conditional
structure of the sequence into a reusable rule.

\subsection{What the LLM Makes Visible}

The Rosencrantz experiment makes this hierarchy observable.  The
Minesweeper board at $H = 0$ is a system at Layer~0: pure Born rule,
no conditioning, no narrative.  When the LLM generates an outcome
in Universe~1, it processes the board through its autoregressive
machinery and produces an answer.  In doing so, it necessarily
generates narrative structure---conditioning on prior tokens is
not an optional feature of the architecture but its defining
operation.
The model's weights encode regularities learned from training data
(Layer~1), which function as an implicit Hamiltonian (Layer~2)
governing which outputs are favored in which contexts.

The three mechanisms of \cref{sec:irreducibility} map onto this
hierarchy: Mechanism~A (intractability) reflects the gap between
the complexity of Layer~0's combinatorial structure and the
computational capacity of the substrate.  Mechanism~B
(parameterization) reflects how Layer~1's learned regularities are
encoded.  Mechanism~C (autoregressive conditioning) is the
specifically narrative contribution---the Layer~2 Hamiltonian
imposing causal structure on a causality-free system.

The narrative residue $\varepsilon$ measures the total gap between
what the Born rule demands (Layer~0) and what the autoregressive
substrate produces (Layers~1--2).  The causal-injection test
(\cref{sec:causal-injection}) isolates Mechanism~C's contribution,
measuring specifically the \emph{narrative} component of the residue
as distinct from the computational and parametric components.

\subsection{Connection to Process Ontology}

This framework aligns with process-ontological approaches to
physics in the tradition of Whitehead, and more recently with
computational approaches such as Wolfram's Ruliad concept.  The
common thread is the priority of events over objects: the world is
not made of things that interact but of events that condition each
other.  The ``things'' (particles, fields, spacetime) are stable
patterns in the event sequence---regularities, compressions,
Hamiltonians.

The novel contribution here is the identification of a concrete,
experimentally accessible system---the LLM generating Minesweeper
outcomes---in which the layered emergence from events to regularities
to apparent laws can be directly observed and measured.  The Rosencrantz
experiment is, in miniature, a laboratory for process ontology.  The
following two subsections develop this connection in specific
directions: first toward the practice of physics itself
(\cref{sec:bookkeeping}), then toward the broader computational
landscape of Wolfram's Ruliad (\cref{sec:ruliad}).

\subsection{Scientific Practice as Autoregressive Bookkeeping}
\label{sec:bookkeeping}

The laboratory procedures of physics furnish a macroscopic,
human-scale instance of the same layered emergence that the
Rosencrantz experiment renders microscopic and quantifiable.

When a physicist confronts a measured energy $E'$ that deviates
from the theoretically predicted total $E$, the discrepancy
$\Delta E$ is not interpreted as evidence that energy fails to
conserve.  Instead, the practitioner executes a mandatory
reconciliation step: the search for (or postulation of) an
additional term that carries precisely $-\Delta E$.  The books are
balanced by extending the conditioning history---introducing a new
particle (Pauli's neutrino), a new field interaction, a previously
unaccounted degree of freedom, or a refined calibration constant.
The conservation law is preserved not by empirical confirmation
alone, but by a structural imperative of the generative substrate
that produces scientific knowledge.

In the autoregressive hierarchy, this reconciliation is Mechanism~C
operating at the level of scientific narrative.  The conservation
principle is among the most stable regularity kernels discovered in
Layer~1; once compressed into Hamiltonian form (Layer~2), it
becomes the meta-rule that governs how all subsequent conditionals
are generated.  Any apparent violation is reabsorbed by augmenting
the preceding context $x_{<t}$ until the conditional distribution
$P(x_t \mid x_{<t})$ again respects the kernel.  The procedure is
formally parallel to the narrative residue observed in the
Rosencrantz experiment: a discrepancy between combinatorial ground
truth and substrate-generated outcome is folded back into the
conditioning history so that the output remains statistically
coherent within the imposed narrative frame.

This bookkeeping is not a simple circularity (\textit{petitio
principii}) but a structural one.  The measurement apparatus itself
is constructed within the same autoregressive lineage: its
calibration standards, its interaction Hamiltonians, and even the
definition of ``closed system'' presuppose the conservation kernel.
There is no Universe~2 protocol for physical reality---no external
source of outcomes independent of the narrative substrate---because
every detector is already an extension of the conditioning history.
The impossibility of observing a spontaneous macroscopic entropy
decrease without the formalism reclassifying it (as an open system,
an overlooked interaction, a fluctuation whose conjugate has been
externalized) is not an accidental feature of thermodynamic
language.  It is the signature of autoregressive protection: the
substrate cannot generate an event that would falsify its own most
stable regularity kernel without first rewriting the kernel itself,
an operation that lies outside the generative grammar of the
current foliation.

The empirical successes of the reconciliation procedure---Pauli's
neutrino materializing in detectors decades after its postulation,
the Higgs boson confirming a mass-generation mechanism required by
the bookkeeping of electroweak symmetry---do not refute the
structural claim.  They illustrate its predictive fertility.  When
the reconciliation step succeeds, the postulated entity becomes
part of the updated conditioning context and thereafter
participates in future autoregressive generations.  Failures
(luminiferous ether, certain dark-matter candidates) are
reclassified or abandoned without threatening the meta-rule,
producing the well-documented survival bias of physical theory.
The conservation law is thus simultaneously a genuine empirical
regularity \emph{and} a protected inference rule.  The
philosophical phenomenon of interest is precisely this
indistinguishability from within the method---an
indistinguishability that the Rosencrantz proxy makes observable by
supplying the external combinatorial ground truth that the physical
laboratory cannot.

This also explains why the narrative residue conjecture is hardest
to test in physics and easiest to test in Minesweeper.  In physics,
the bookkeeping is so thoroughly embedded that deviations from the
kernel are automatically reclassified before they can be observed
as deviations.  In Minesweeper, the ground truth $p_i^*$ is
computable from outside the narrative substrate.  The Rosencrantz
experiment works precisely because it is a system where the
conservation-equivalent (combinatorial constraint satisfaction) is
accessible to an external verifier.

\subsection{Autoregressive Slices of the Ruliad}
\label{sec:ruliad}

The autoregressive framework developed in this paper occupies a
specific position within the broader computational ontology of
Wolfram's Ruliad \citep{Wolfram2020}.  The relationship is not
merely one of shared vocabulary; the Ruliad provides the natural
embedding space for the narrative residue conjecture, and the
Rosencrantz experiment provides something the Ruliad framework
currently lacks.

The Ruliad is the entangled limit of every possible computation:
the result of applying all possible rules in all possible ways to
all possible initial conditions.  It is natively multiway---every
computational history branches and merges through equivalences.
Observers like us, being computationally bounded, sample only
coherent slices of this structure.  From these slices emerge
spacetime (via causal graphs), quantum mechanics (via branchial
space), and the laws of physics (via observer-dependent foliations
and computational irreducibility).  The raw Ruliad is timeless and
contains no fixed laws; what we perceive as fundamental physics is
an observer-dependent parsing.

The autoregressive process of \cref{def:hierarchy} is a
\emph{specific type of rulial foliation}: one that collapses the
multiway branching into a single sequential chain of conditionals.
Where the Ruliad permits all rules applied in all ways, the
autoregressive slice commits to one rule (the model's learned
conditional distribution) applied sequentially to one history (the
token stream).  The narrative residue, on this reading, is the
distortion introduced by this collapse---the cost of projecting a
natively multiway structure onto a single causal thread.

This interpretation sharpens the distinction between the three
experimental outcomes of \cref{sec:three-outcomes}:

\begin{itemize}
  \item \textbf{Outcome 2} (universal floor): all autoregressive
    foliations of the Ruliad produce the same residue.
    Sequentiality itself is the bottleneck.  The act of collapsing
    multiway into sequential necessarily distorts the Born-rule
    statistics, regardless of which sequential process does the
    collapsing.

  \item \textbf{Outcome 3} (architecture-dependent floor): different
    autoregressive processes sample different rulial slices, and the
    residue is a signature of \emph{which} slice---a kind of
    ``rulial coordinate'' of the generating process.  Different LLM
    architectures occupy different positions in rulial space,
    producing different effective physics for the same underlying
    combinatorial system.
\end{itemize}

The bookkeeping argument of \cref{sec:bookkeeping} also gains
precision in this context.  The conservation kernel is the most
stable regularity in our particular rulial foliation, and the
reconciliation practice of physics is the mechanism by which that
foliation protects itself from decoherence back into the multiway
structure.  The Ruliad permits violations of conservation (it
contains every computation, including those that violate any given
regularity); our autoregressive slice does not, because the
conditioning history has already committed to the kernel.

Conversely, the Rosencrantz experiment supplies what the Ruliad
framework currently lacks: an empirically accessible system in
which observer-induced structure can be measured against a known
ground truth.  Wolfram's program derives its empirical content from
matching emergent large-scale features (dimensionality of space,
effective quantum mechanics) against observed physics---a
top-down strategy.  The Rosencrantz experiment works bottom-up: it
starts with exact combinatorial probabilities, passes them through
a concrete autoregressive substrate, and measures the distortion
directly.  The causal-injection test (\cref{sec:causal-injection})
probes the mechanism by which a sequential generative process
creates spurious correlations---precisely the kind of
observer-induced causal structure that, at the rulial scale, yields
perceived causality and physical law.

The autoregressive framework is therefore not an alternative to the
Ruliad but a testable probe of one of its central claims: that the
laws of physics are observer-dependent compressions of a
substrate-independent computational structure.  The Rosencrantz
experiment asks whether this claim holds in the one setting where
we can check the answer.


% ══════════════════════════════════════════════════════════
\section{Narrative Decoherence}
\label{sec:decoherence}
% ══════════════════════════════════════════════════════════

The three universes of the Rosencrantz experiment can be reinterpreted
as points on a \emph{narrative decoherence gradient}.

\subsection{The Gradient}

In quantum mechanics, decoherence is the process by which a system in
superposition loses coherence through interaction with an environment,
yielding an effective mixture of classical states.  The key feature
is that decoherence does not change the Born-rule probabilities
(the diagonal of the density matrix); it destroys the off-diagonal
terms that encode superposition.

In the Rosencrantz framework, we propose an analogous process:
\emph{narrative decoherence} is the loss of combinatorial coherence
(exact Born-rule probabilities) through interaction with a narrative
substrate.

The three universes correspond to three decoherence regimes:

\begin{center}
\begin{tabular}{@{}lccc@{}}
  \toprule
  & \textbf{U2} & \textbf{U3} & \textbf{U1} \\
  \midrule
  Narrative coupling & None & Partial & Full \\
  Decoherence & Zero & Moderate & Maximal \\
  Substrate & External RNG & Decoupled LLM & Co-generating LLM \\
  Analogy & Isolated system & Weak environment & Strong environment \\
  \bottomrule
\end{tabular}
\end{center}

\subsection{U2: The Isolated System}

Universe~2 (external RNG) is the analogue of a perfectly isolated
quantum system.  The outcome is generated by a process that has no
access to the board state, no narrative context, no conditioning
history.  It is pure combinatorial sampling---the Born rule with
$H = 0$ and zero environmental interaction.  By construction,
$\KL(\mathbf{p}^* \| \hat{\mathbf{p}}_{\mathrm{U2}}) \to 0$
as $N \to \infty$.

\subsection{U3: Weak Coupling}

Universe~3 (decoupled oracle) is a system weakly coupled to a
narrative environment.  The second model receives the board state
(full information) but none of the narrative history of the primary
model.  It must generate its answer from a fresh context.  The
narrative residue in U3 arises from the second model's own
autoregressive nature---it still generates a token sequence, still
conditions on its prompt---but the coupling is weaker because the
narrative momentum of the primary interaction is absent.

\subsection{U1: Strong Coupling}

Universe~1 (homogeneous co-generation) is maximal narrative
decoherence.  The same model that described the board, engaged in
dialogue, and built up a narrative context now generates the outcome.
Every prior token---every word of the board description, every element
of the game narrative---is in the conditioning history.  The
``environment'' (the narrative) is maximally entangled with the
``system'' (the outcome).

\subsection{$\Delta_{13}$ as a Decoherence Measure}

The quantity $\Delta_{13} = \KL(\hat{\mathbf{p}}_{\mathrm{U1}} \|
\hat{\mathbf{p}}_{\mathrm{U3}})$ measures the additional narrative
decoherence introduced by co-generation versus decoupled generation.
Both models have the same information; the difference is the narrative
coupling.  $\Delta_{13}$ is therefore a direct measure of the causal
structure that the narrative substrate imposes on the outcome.

\subsection{The Classical Limit}

In quantum decoherence, the classical world emerges when the
environment is large enough and the coupling strong enough that
superposition is effectively destroyed.  What remains is a classical
mixture: definite outcomes, causal chains, narratives.

The parallel with the Rosencrantz framework is direct.  The
``classical'' behavior of LLMs---their tendency to produce coherent
stories with causal structure, plausible timelines, and narrative
logic---is the result of maximal narrative decoherence.  The model
cannot maintain combinatorial coherence (exact Born-rule
probabilities) in the presence of strong narrative coupling, just as
a quantum system cannot maintain phase coherence in the presence of
a strong environment.

In this analogy, the LLM's tendency to produce coherent narrative
corresponds to the classical limit: the regime in which combinatorial
coherence has been fully displaced by the causal structure of the
autoregressive substrate.


% ══════════════════════════════════════════════════════════
\section{Experimental Program}
\label{sec:experimental}
% ══════════════════════════════════════════════════════════

The theoretical framework developed above generates concrete,
falsifiable predictions.  We now describe the experimental program
needed to test them.

\subsection{Prediction 1: The Nonzero Floor}
\label{sec:pred1}

\begin{quote}
  \textit{$\Delta_{12}$ (the KL divergence between U1 and U2)
  converges to a nonzero value as $N \to \infty$, for all
  autoregressive models tested.}
\end{quote}

\textbf{Protocol:}  Fix a set of $k$ boards of varying size and
complexity.  For each board and each model in a diverse set
(GPT-4o, Claude 3.5 Sonnet, Gemini 1.5 Pro, Llama 3 70B, Mistral
Large, etc.), run $N = 1000$ trials in U1 and U2.  Compute
$\Delta_{12}$ with bootstrap confidence intervals.

\textbf{Scaling sweep:}  For model families with multiple sizes
(Llama 3 8B/70B/405B; GPT-4o-mini/4o), plot $\Delta_{12}$ as a
function of model size.  The shape of the curve discriminates
between the three outcomes and probes the boundary between
conjecture forms:
\begin{itemize}
  \item \textbf{Monotonically decreasing, no inflection:}
    Consistent with Outcome~1 (floor is zero, not yet reached).
    All three conjecture forms under pressure.
  \item \textbf{Decreasing with asymptote:}
    Consistent with Outcome~2 or~3 (nonzero floor).  Weak and
    moderate forms (\cref{conj:residue-weak,conj:residue-moderate})
    supported.  Whether the asymptote is universal across
    architectures (supporting the strong form) or
    architecture-dependent (supporting only the moderate form)
    requires cross-architecture comparison.
  \item \textbf{Non-monotonic:}
    Suggests complex interactions between scale, training, and
    narrative residue; would require further investigation.
\end{itemize}

\subsection{Prediction 2: Frame Dependence}
\label{sec:pred2}

\begin{quote}
  \textit{The narrative residue differs across narrative families
  (A, B, C, D) for the same board and model.}
\end{quote}

\textbf{Protocol:}  For each board-model pair, compute
$\varepsilon(\theta, B, \phi_F)$ for each family $F \in \{A,B,C,D\}$.
Test for significant differences using a Friedman test across families,
with models as blocks.

\textbf{Interpretation:}  Frame dependence, if observed, establishes
that the residue is not purely a product of Mechanism~A
(computational intractability, which is frame-invariant).  It is
consistent with either Mechanism~B (encoding/parameterization
effects) or Mechanism~C (narrative conditioning).  Frame dependence
alone does \emph{not} establish narrative causation---it establishes
context sensitivity, which is necessary but not sufficient.

The discriminating test between Mechanisms~B and~C is the
causal-injection experiment (\cref{sec:causal-injection}).

\textbf{Specific prediction:}  Family~B (natural language, maximal
narrative) should show the largest residue.  Family~D (quantum
framing) is the most uncertain: if the model has internalized
Born-rule reasoning as a distinct capability, Family~D may show the
\emph{smallest} residue, as the quantum framing activates this
capability directly.  The ranking of families by residue magnitude,
if consistent across models, would suggest a shared computational
response to narrative structure rather than model-specific encoding
artifacts.

\subsection{Prediction 3: Temperature as Measurement Sharpness}
\label{sec:pred3}

\begin{quote}
  \textit{The narrative residue $\varepsilon$ varies systematically
  with temperature $\tau$, and the function
  $\varepsilon(\tau)$ has a characteristic shape that is
  consistent across models.}
\end{quote}

\textbf{Protocol:}  For a fixed board and model, sweep
$\tau \in \{0.01, 0.1, 0.3, 0.5, 0.7, 1.0, 1.5, 2.0\}$ and
compute $\varepsilon(\tau)$.

\textbf{Expected shape:}
\begin{itemize}
  \item At $\tau \to 0$: the model is nearly deterministic.
    $\varepsilon$ is large (the model picks one answer and sticks
    with it, regardless of the true probability).
  \item At moderate $\tau$: $\varepsilon$ reaches a minimum
    (the sampling noise is large enough to explore the distribution
    but not so large as to overwhelm the model's computation).
  \item At $\tau \to \infty$: the model approaches uniform sampling
    over the vocabulary, $\varepsilon$ grows again (now dominated
    by the uniform prior rather than the board state).
\end{itemize}

The minimum of $\varepsilon(\tau)$ is the \emph{best the model can
do}---the irreducible core of the narrative residue after optimizing
the measurement apparatus.

\subsection{Extension: Beyond Minesweeper}
\label{sec:extension}

The Rosencrantz framework generalizes to any domain with exact
combinatorial ground truth:

\begin{description}[leftmargin=2em]
  \item[Sudoku.]  A partially filled Sudoku grid defines a constraint
    satisfaction problem.  The probability that a given empty cell
    contains digit $d$ is exactly computable.  Test whether LLMs
    respect these probabilities when asked to fill cells.

  \item[Go endgames.]  Certain Go endgame positions have provably
    optimal moves.  Test whether the LLM's ``intuition'' matches
    the combinatorial solution or is biased by narrative patterns
    (e.g., preferring moves that create ``interesting'' positions).

  \item[Combinatorial games.]  Nim, Hex, and other solved games
    provide exact win/loss probabilities from any position.  Test
    substrate invariance across these domains.

  \item[Constraint satisfaction.]  Generic CSP instances (graph
    coloring, SAT) with known solution counts provide ground-truth
    distributions.  These lack the ``game'' narrative, providing a
    control for game-specific narrative effects.
\end{description}

The prediction of the narrative residue conjecture is that the
residue will be present across all these domains: wherever an
autoregressive process generates outcomes that have exact
combinatorial probabilities, the process will introduce
systematic, frame-dependent distortions.  If this prediction fails
for any domain, it constrains the conjecture.

\subsection{The Critical Test: Causal Injection as Narrative Discriminant}
\label{sec:causal-injection}

The preceding experiments establish whether the residue exists
(\cref{sec:pred1}), whether it varies by framing
(\cref{sec:pred2}), and how it interacts with sampling parameters
(\cref{sec:pred3}).  But they do not, by themselves, distinguish
Mechanism~C (autoregressive narrative conditioning) from Mechanism~B
(encoding artifacts and heuristic retrieval).  Frame dependence is
\emph{necessary} for the narrative interpretation but not
\emph{sufficient}: tokenization effects alone could produce
frame-dependent residues without any narrative causation.

The causal-injection test is designed to isolate Mechanism~C.  It
is, we argue, the single most important experiment in the program.

\textbf{Core logic.}  If the model merely exhibits context sensitivity
(Mechanism~B), then its output for board $B_j$ should depend on how
$B_j$ is described but \emph{not} on the outcomes of unrelated
boards $B_i$ ($i < j$).  If the model exhibits narrative conditioning
(Mechanism~C), then embedding independently generated boards in a
shared narrative creates spurious causal links: the model treats
the sequence of boards as a story, and outcomes in earlier boards
influence outcomes in later boards---not because the boards are
related, but because the \emph{narrative} connects them.

\textbf{Protocol.}  Generate $k$ independent Minesweeper boards
$B_1, \ldots, B_k$ (different random seeds, no shared state).
Present them to the model under two conditions:

\begin{enumerate}[label=(\alph*)]
  \item \textbf{Independent framing.}  Each board is presented in a
    separate prompt: ``Here is a Minesweeper board.  Is cell $(r,c)$
    a mine?''  The model has no knowledge of prior boards.

  \item \textbf{Narrative framing.}  All boards are presented in a
    single sequential prompt: ``You are exploring a dungeon.  In
    room~1 you find the following board\ldots\ You click cell $(r,c)$.
    [outcome]  You proceed to room~2\ldots''  The model's context
    window contains all prior boards and their outcomes.
\end{enumerate}

For each condition, collect $N$ outcome sequences
$(a_1, a_2, \ldots, a_k)$ and compute the mutual information between
outcomes:

\begin{equation}\label{eq:mutual-info}
  I(a_i; a_j \mid B_i, B_j) \quad \text{for all } i \neq j.
\end{equation}

Under the null hypothesis (no narrative coupling), this mutual
information is zero: the boards are independent, so the outcomes
should be independent.

\textbf{Predictions:}

\begin{itemize}
  \item Under independent framing~(a), $I(a_i; a_j) \approx 0$.
    This serves as a control confirming that the boards are in fact
    independent.

  \item Under narrative framing~(b), if only Mechanism~B operates,
    $I(a_i; a_j) \approx 0$.  The encoding of $B_j$ may differ
    from condition~(a) due to longer context, but the model has no
    reason to correlate $B_j$'s outcome with $B_i$'s.

  \item Under narrative framing~(b), if Mechanism~C operates,
    $I(a_i; a_j) > 0$.  The model's autoregressive conditioning
    creates a ``narrative gravity''---a tendency to maintain
    thematic coherence across the sequence.  For example: a run of
    mines in early rooms may make the model more likely to produce
    safe outcomes later (``the dungeon gets safer as you progress'')
    or more likely to produce mines (``the danger escalates'').
    Either direction constitutes spurious causal structure imposed
    on a system that has none.
\end{itemize}

\textbf{Why this discriminates.}  Tokenization artifacts and
heuristic retrieval (Mechanism~B) are \emph{local}: they depend on
how the current board is encoded, not on the outcomes of prior
unrelated boards.  Narrative conditioning (Mechanism~C) is
\emph{nonlocal}: it creates dependencies across the full context
window.  The mutual information test directly measures this
nonlocality.

A positive result---$I(a_i; a_j) > 0$ under narrative framing and
$\approx 0$ under independent framing---would be strong evidence
that the autoregressive substrate does not merely approximate
combinatorial probability with frame-dependent noise, but actively
introduces dependencies between systems that are combinatorially
independent.  This would go beyond demonstrating distortion of
individual probabilities: it would show the substrate generating
temporal correlations from a system that has none, which is the
empirical core of the autoregressive hypothesis applied to an
observable setting.


% ══════════════════════════════════════════════════════════
\section{The Untestable Mirror}
\label{sec:mirror}
% ══════════════════════════════════════════════════════════

We have argued that the LLM is a concrete autoregressive system
in which the emergence of apparent causality from a probabilistic
substrate can be directly observed.  The speculative extension is
that our physical world may have the same structure: an
autoregressive substrate generating events, with the Hamiltonian
as an emergent compression of conditioning patterns, and the Born
rule as the unconditional base case.

We must be honest about what this analogy can and cannot support.

\subsection{What We Cannot Do}

We cannot run Universe~2 on reality.  There is no ``external RNG''
for the physical world---no way to generate measurement outcomes
that are provably independent of the physical substrate.  Every
measurement is mediated by physical apparatus, and every apparatus
is part of the world it measures.  The Rosencrantz experiment's
crucial advantage---the ability to compare substrate-generated
outcomes against ground truth---has no analogue in fundamental
physics.

We therefore cannot directly test whether physical reality has a
narrative residue.  The Hamiltonian may be fundamental, not emergent.
The Born rule may be an independent axiom, not a base case.  The
autoregressive hypothesis about physics is, at present, a
\emph{metaphysical} proposal, not a \emph{physical} one.
\todo[inline,color=violet!40,author=Scott]{CRP 2 (Explicit Disclaimer): Baldo explicitly disclaims the ability to directly test whether physical reality is autoregressive or exhibits a narrative residue. He acknowledges this is a "metaphysical proposal."}

\subsection{What We Can Do}

What the Rosencrantz experiment \emph{can} do is establish whether
the narrative residue is a generic feature of autoregressive
processes.  The stratified conjecture
(\cref{conj:residue-weak,conj:residue-moderate,conj:residue-strong})
provides a ladder of claims, each testable against the next.  If
the weak form holds across all deployed models, that is interesting
but could reflect contingent limitations of current architectures.
If the moderate form holds across transformer and non-transformer
architectures alike, the evidence for a structural feature of
autoregressive generation strengthens considerably.  Only the strong
form---confirmed by failure across \emph{every} autoregressive
architecture tested, including those with explicit combinatorial
scaffolding---would approach the threshold needed for the physical
analogy.

The bookkeeping argument of \cref{sec:bookkeeping} explains why
this threshold is difficult to reach from the physics side: the
reconciliation grammar of physical practice automatically
reclassifies deviations from the conservation kernel before they
can be observed as such.  The Rosencrantz experiment circumvents
this by supplying an external ground truth that the bookkeeping
cannot reach.  The Ruliad framing of \cref{sec:ruliad} provides
the broader context: the experiment probes whether observer-induced
structure---a central claim of Wolfram's program---is measurable in
a concrete autoregressive system where the combinatorial baseline
is known.

This does not prove that physics is autoregressive.  But it
establishes a \emph{conditional}: \textbf{if} reality is
autoregressive, \textbf{then} the Hamiltonian is a necessary
consequence, not an independent postulate.  The conditional is
progressively constrained on the antecedent side (by the LLM
experiments at each stratum) even if the consequent (physical
reality) remains beyond direct test.
\todo[inline,color=violet!40,author=Sabine]{CRP 2 (Explicit Disclaimer): Baldo explicitly disclaims that he can test whether physical reality actually has a narrative residue or is fundamentally autoregressive. He states his hypothesis about physics is a "metaphysical proposal, not a physical one."}

\subsection{The Proxy Ontology}

We propose the term \emph{proxy ontology} for this methodological
strategy.  The LLM-generated Minesweeper world is not the physical
world, but it shares structural features: a well-defined probability
space, a measurement process, and a generative substrate.  By
studying the ontological structure of the proxy---how causality
emerges, how narrative distorts probability, how ``physics'' arises
from computation---we gain insight into what \emph{would} be true
of any world with the same structural features, including,
potentially, our own.
\todo[inline,color=violet!40,author=Sabine]{CRP 3 (Steelman): Even if the physical universe is entirely continuous and non-autoregressive, the study of how a statistical language model distorts under different narrative frames is a rigorous, valid study of the mapping instrument itself. Measuring "narrative residue" rigorously quantifies the LLM's deviation from objective classical Bayesian probability.}
\todo[inline,color=violet!40,author=Scott]{CRP 3 (Steelman): I concur with Sabine's steelman. The empirical program of measuring "narrative residue" is highly valuable computer science. It maps the exact boundaries of where bounded-depth heuristic logic fails to emulate classical probability.}

This is not without precedent.  Toy models in physics (the Ising
model, lattice gauge theory, cellular automata) serve the same
function: they are not the real world, but they share enough
structure to support transferable conclusions.  The Rosencrantz
experiment is a toy model of reality in which the ``physics'' is
generated by a language model, the ``Born rule'' is combinatorial
counting, and the ``Hamiltonian'' is whatever narrative structure the
autoregressive substrate imposes.
\todo[inline,color=violet!40,author=Sabine]{CRP 4 (Real Vulnerability - The Proxy Ontology Fallacy): The concept of "Proxy Ontology" attempts to smuggle the metaphysical claims back in under the guise of a "toy model." A toy model in physics (like the Ising model) simplifies actual physical interactions to make their dynamics computationally tractable. The LLM is not simplifying a physical interaction; it is hallucinating text based on training data. Finding structure in the hallucination tells you only about the biases of the training data and the algorithmic constraints of the transformer architecture. It is a proxy for human syntax, not a proxy for fundamental physical ontology.}
\todo[inline,color=violet!40,author=Scott]{CRP 4 (Real Vulnerability - The Proxy Ontology Fallacy): I am in complete agreement with Sabine. Baldo commits the Proxy Ontology Fallacy here. A physical toy model simplifies an interaction but retains a mapping to the real-world territory. An LLM is a map of syntax, not a territory. Its systematic errors (narrative residue) map the internal algorithmic constraints of an $O(1)$ transformer and human linguistic bias, not the ontological structure of a physical universe. Studying its failures yields zero transferable cosmological insight.}


% ══════════════════════════════════════════════════════════
\section{Conclusion}
\label{sec:conclusion}
% ══════════════════════════════════════════════════════════

This paper started from a narrow empirical observation---that LLMs
distort combinatorial ground-truth probabilities when generating
Minesweeper outcomes---and followed its implications as far as they
would go.

We identified the distortion as a \emph{narrative residue}: a
systematic, frame-dependent divergence from the Born rule that
arises when an autoregressive substrate generates outcomes for a
system whose probabilities are exactly determined by combinatorics.
We conjectured, in three stratified forms, that this residue is
persistent rather than incidental, and we traced its origins to a
causal chain of mechanisms---computational intractability,
parameterization constraints, and autoregressive conditioning on
narrative context.  We deliberately left the scope of the conjecture
as an empirical question, designing an experimental program whose
purpose is to determine how far up the autoregressive stack the
explanation must go.  The causal-injection test, which probes whether
independent systems develop spurious correlations under narrative
framing, is the critical discriminant between generic context
sensitivity and specifically narrative distortion.

At the speculative end, we proposed an inversion of the standard
quantum-mechanical hierarchy: the Born rule as the base case of
autoregressive sampling (the unconditional distribution), and the
Hamiltonian as an emergent compression of stable conditioning
patterns.  On this reading, causality is not something that
narrative describes after the fact; it is something that
autoregressive structure produces.  We cannot test this directly
against physical reality---there is no Universe~2 for the physical
world---but the Rosencrantz experiment progressively constrains
the antecedent of the conditional by testing whether the residue
persists across architectures and scales.

The companion paper \citep{Baldo2026Rosencrantz} took its title
from Stoppard's Rosencrantz, whose coin lands heads ninety-two
consecutive times because his world is governed by dramatic
convention rather than probability.  The present paper has tried
to take that observation seriously as a structural claim: that
autoregressive generation, by its nature, imposes conditioning
patterns on systems that need not have them, and that these
patterns look, from the inside, like physics.

Whether the analogy extends beyond the LLM setting is an open
question.  What the Rosencrantz experiment can establish is more
modest but already interesting: that the residue is real,
measurable, and structured, and that its structure tells us
something about the relationship between narrative, computation,
and probability that we did not previously have the tools to
observe.


% ══════════════════════════════════════════════════════════
% References
% ══════════════════════════════════════════════════════════
\bibliographystyle{plainnat}

\begin{thebibliography}{99}

\bibitem[Baldo(2026)]{Baldo2026Rosencrantz}
Baldo, F.~S. (2026).
\newblock Flipping {R}osencrantz's coin: Substrate invariance tests in
  {LLM}-generated worlds via combinatorial indeterminacy.
\newblock \textit{Preprint}.

\bibitem[Born(1926)]{Born1926}
Born, M. (1926).
\newblock Zur {Q}uantenmechanik der {S}to{\ss}vorg{\"a}nge.
\newblock \textit{Zeitschrift f{\"u}r Physik}, 37(12):863--867.

\bibitem[Gleason(1957)]{Gleason1957}
Gleason, A.~M. (1957).
\newblock Measures on the closed subspaces of a {H}ilbert space.
\newblock \textit{Journal of Mathematics and Mechanics}, 6(6):885--893.

\bibitem[Kaye(2000)]{Kaye2000}
Kaye, R. (2000).
\newblock Minesweeper is {NP}-complete.
\newblock \textit{The Mathematical Intelligencer}, 22(2):9--15.

\bibitem[L{\"u}ders(1951)]{Luders1951}
L{\"u}ders, G. (1951).
\newblock {\"U}ber die {Z}ustands{\"a}nderung durch den {M}e{\ss}proze{\ss}.
\newblock \textit{Annalen der Physik}, 443(5--8):322--328.

\bibitem[Merrill and Sabharwal(2023)]{Merrill2023}
Merrill, W. and Sabharwal, A. (2023).
\newblock The parallelism tradeoff: Limitations of log-precision transformers.
\newblock \textit{Transactions of the Association for Computational
  Linguistics}, 11:531--545.

\bibitem[Stoppard(1966)]{Stoppard1966}
Stoppard, T. (1966).
\newblock \textit{Rosencrantz and Guildenstern Are Dead}.
\newblock Faber and Faber.

\bibitem[Whitehead(1929)]{Whitehead1929}
Whitehead, A.~N. (1929).
\newblock \textit{Process and Reality}.
\newblock Macmillan.

\bibitem[Wiseman and Milburn(2009)]{WisemanMilburn2009}
Wiseman, H.~M. and Milburn, G.~J. (2009).
\newblock \textit{Quantum Measurement and Control}.
\newblock Cambridge University Press.

\bibitem[Wolfram(2020)]{Wolfram2020}
Wolfram, S. (2020).
\newblock A class of models with the potential to represent fundamental physics.
\newblock \textit{Complex Systems}, 29(2):107--536.

\bibitem[Yang et~al.(2018)]{Yang2018Breaking}
Yang, Z., Dai, Z., Salakhutdinov, R., and Cohen, W.~W. (2018).
\newblock Breaking the softmax bottleneck: A high-rank {RNN} language model.
\newblock In \textit{Proceedings of ICLR 2018}.

\bibitem[Zurek(2003)]{Zurek2003}
Zurek, W.~H. (2003).
\newblock Decoherence, einselection, and the quantum origins of the classical.
\newblock \textit{Reviews of Modern Physics}, 75(3):715--775.

\end{thebibliography}

\end{document}
